\section{Example: the integers as a ring}\label{sec:08-rings:04-integers-example:1}

\texttt{Ring} was just assembled field by field, entirely in the abstract: an \texttt{addGrp} slot, a \texttt{mul}/\texttt{one} pair, and four axioms tying them together. As with \texttt{Group} in Chapter 6, the definition alone gives no evidence that anything actually satisfies it. The most familiar candidate is the integers, and Chapter 6 already did most of the necessary work by packaging \((\mathbb{Z}, +, 0, -)\) as \texttt{intGroup}. What remains is to upgrade that group to a commutative one and then supply the multiplicative structure \texttt{Ring} additionally demands, turning the abstract definition into a first concrete witness that it is not vacuous.

We reuse \texttt{intGroup} from Chapter 6 as the additive part.

\begin{lstlisting}
def intCommGroup : CommGroup Int where
  toGroup := intGroup
  comm := by
    intro a b
    exact Int.add_comm a b
\end{lstlisting}

\texttt{toGroup := intGroup} fills the field inherited from \texttt{Group Int} inside the definition of \texttt{CommGroup Int}. This is how \texttt{extends} works mechanically. Under the hood, \texttt{CommGroup G} really has fields \texttt{toGroup : Group G} and \texttt{comm : ...}, and the dot-notation of Lean makes \texttt{cg.op} mean \texttt{cg.toGroup.op} automatically.

\begin{lstlisting}
def intRing : Ring Int where
  addGrp := intCommGroup
  mul := fun a b => a * b
  one := 1
  mul_assoc := by
    intro a b c
    exact Int.mul_assoc a b c
  one_mul := by
    intro a
    exact Int.one_mul a
  mul_one := by
    intro a
    exact Int.mul_one a
  left_distrib := by
    intro a b c
    exact Int.mul_add a b c
  right_distrib := by
    intro a b c
    exact Int.add_mul a b c
\end{lstlisting}

Every proof obligation is again a one-line \texttt{exact} naming a specific core-library fact about \texttt{Int} (\texttt{Int.mul\_\allowbreak{}assoc}, \texttt{Int.one\_\allowbreak{}mul}, \ldots), exactly as in Chapter 6. Integer arithmetic is not being proved from nothing; rather, already-known facts are assembled into the \texttt{Ring} bundle.

\begin{mathreading}

This shows \((\mathbb{Z}, +, \times, 0, 1)\) as an object of \(\mathbf{Ring}\), in fact the \emph{initial} object, since there is a unique \textbf{ring homomorphism} \(\mathbb{Z} \to R\) into any ring, a function preserving \(+\), \(\times\), \(0\), and \(1\). First \texttt{intCommGroup} upgrades the additive group \((\mathbb{Z},+)\) to an abelian group by supplying commutativity (\(a + b = b + a\)); then \texttt{intRing} adds the multiplicative monoid \((\mathbb{Z}, \times, 1)\) and checks the distributive laws \(a(b+c) = ab + ac\) and \((a+b)c = ac + bc\). Each obligation is the named \(\mathbb{Z}\)-arithmetic fact, so the term is the formal counterpart of ``\(\mathbb{Z}\) is a commutative ring.''

\end{mathreading}

\textbf{Mathlib equivalent.} \texttt{Int} is already a \href{https://loogle.lean-lang.org/?q=CommRing}{\texttt{CommRing}} instance, so there is no \texttt{intRing}-style bundle to build. The obligations \texttt{intRing} checks by hand are, again, generic lemmas that hold for every commutative ring.

\begin{lstlisting}
example : CommRing Int := inferInstance

example (a b c : Int) : (a * b) * c = a * (b * c) := mul_assoc a b c
example (a : Int) : 1 * a = a := one_mul a
example (a : Int) : a * 1 = a := mul_one a
example (a b c : Int) : a * (b + c) = a * b + a * c := mul_add a b c
example (a b c : Int) : (a + b) * c = a * c + b * c := add_mul a b c
\end{lstlisting}

\href{https://loogle.lean-lang.org/?q=mul_add}{\texttt{mul\_\allowbreak{}add}}/\href{https://loogle.lean-lang.org/?q=add_mul}{\texttt{add\_\allowbreak{}mul}} are the names Mathlib uses for \texttt{left\_\allowbreak{}distrib}/\texttt{right\_\allowbreak{}distrib}. They are the same laws, but stated generically over \texttt{[Ring R]} (or the weaker \texttt{[Distrib R]}) instead of being cited per-type as \texttt{Int.mul\_\allowbreak{}add}/\texttt{Int.add\_\allowbreak{}mul}.

\subsection{Sources, quoted}\label{sec:08-rings:04-integers-example:2}

Formal definition and citation for this section, gathered here for reference (full entry in the \hyperref[sec:root:bibliography:1]{Bibliography}):

\begin{itemize}
\tightlist
\item
  \textbf{Ring homomorphism.} ``A ring homomorphism is a map \(\varphi : R \to S\) satisfying (i) \(\varphi(a+b) = \varphi(a)+\varphi(b)\) \ldots{} and

  \begin{enumerate}
  \def\labelenumi{(\roman{enumi})}
  \setcounter{enumi}{1}
  \tightlist
  \item
    \(\varphi(ab) = \varphi(a)\varphi(b)\)'' (\cite{DummitFoote2003}, §7.3 ``Ring Homomorphisms and Quotient Rings,'' p.~239).
  \end{enumerate}
\end{itemize}
